\label{sec:implementation}

The repository intentionally keeps three implementation strata.

\paragraph{Pure JAX reference paths.}
These are the most compact differentiable implementations in the codebase.
They are useful for formulation checks, serial references, and debugging AD
logic. They are not intended to be the primary production route for the largest
distributed benchmarks. Their main scientific role in this paper is to show
that the energy formulations used in the maintained stack can also be expressed
in a compact differentiable form without sparse distributed infrastructure.

\paragraph{FEniCS reference paths.}
The \fenics{} implementations remain valuable on families where they exist,
especially $p$-Laplace, Ginzburg--Landau, and hyperelasticity. They provide a
trusted high-level FEM reference and help distinguish modeling issues from
solver-stack issues. In this repository they are used mainly as
comparison baselines rather than as the mainline production path.

\paragraph{Maintained \jaxpetsc{} path.}
This is the central implementation of the paper. It combines JAX-based local
physics with PETSc distributed vectors, matrices, Krylov solvers, and
preconditioners. Problem-specific kernels stay in JAX; sparse distributed
assembly, nonlinear globalization, and linear algebra stay in PETSc.

\subsection{Autodiff modes}

The maintained stack supports several second-order modes. The precise available
choices depend on the benchmark family, but the common logic is:
\begin{itemize}
  \item \textbf{element AD}: exact local Hessians from differentiating the full
        element energy;
  \item \textbf{constitutive AD}: exact local constitutive tangents assembled
        from quadrature-point energy densities;
  \item \textbf{local SFD / colored SFD}: sparse approximate Hessian recovery
        when exact second-order information is not the appropriate cost tradeoff.
\end{itemize}

This split is scientifically important. It lets us compare exact second-order
information against carefully structured sparse approximations without changing
the rest of the nonlinear solver stack. In other words, the paper is not only a
comparison of frameworks, but also a comparison of derivative strategies inside
one common solver infrastructure.

\begin{figure}[H]
  \centering
  \PaperFrameworkOverviewDiagram
  \caption{Software structure of the manuscript. The maintained \jaxpetsc{}
  implementation is the main production path; \purejax{} and \fenics{} remain
  reference paths where they already exist in the repository.}
  \label{fig:framework-overview}
\end{figure}

\begin{figure}[H]
  \centering
  \PaperDerivativePathDiagram
  \caption{Derivative pathways from scalar energies and quadrature-point
  potentials to distributed PETSc nonlinear solves.}
  \label{fig:derivative-pathways}
\end{figure}

\begin{figure}[H]
  \centering
  \PaperAutodiffModesDiagram
  \caption{Second-order derivative modes exposed in the maintained stack. The
  important engineering choice is that the source of truth remains an energy or
  energy-density expression even when the tangent path changes from element AD
  to constitutive AD or colored SFD.}
  \label{fig:derivative-stack}
\end{figure}

\subsection{Linear solvers and preconditioners}

The maintained families use different Krylov and preconditioner combinations
because the linear systems are qualitatively different:
\begin{itemize}
  \item \textbf{$p$-Laplace:} mostly \code{cg + hypre}, with exact element
        Hessians or local-SFD Hessians in the \jaxpetsc{} path.
  \item \textbf{Ginzburg--Landau:} \code{gmres + hypre}, reflecting the more
        indefinite linearization.
  \item \textbf{Hyperelasticity:} trust-region outer iterations with
        \code{gamg}-preconditioned Krylov solves.
  \item \textbf{Plasticity (2D):} deep-tail multigrid hierarchies whose
        bottlenecks are dominated by the top smoother and repeated Krylov
        work.
  \item \textbf{Plasticity3D:} local PMG on same-mesh
        \code{P4 -> P2 -> P1}, coarse \code{cg + hypre}, and an elastic initial
        guess on the same hierarchy.
  \item \textbf{Topology:} \code{fgmres + gamg} for the mechanics step, coupled
        to a separate distributed design update.
\end{itemize}

The comparison material later in the manuscript also includes a source-repo PMG-shell
variant and deflated GMRES. Those are not the primary maintained defaults, but
they are informative because they reveal how much of the final performance
difference comes from derivative paths, preconditioners, or outer solver logic.

\subsection{Distributed assembly and coloring}

The \jaxpetsc{} implementation uses PETSc-owned ranges, local-plus-ghost
layouts, and owned-row sparse assembly. This matters for both exact and
approximate second-order information. Exact element or constitutive assembly can
write directly into local sparse buffers. Colored SFD requires an additional
graph-coloring step, but once the coloring is known the probe evaluations can
still be carried out in a rank-local manner with only scalar reductions or
owned sparse insertion in the hot loop.

\begin{figure}[H]
  \centering
  \PaperGlobalizationDiagram
  \caption{Globalization choices used across the repository, from pure line
  search to trust-region and hybrid strategies.}
  \label{fig:globalization}
\end{figure}

\begin{figure}[H]
  \centering
  \PaperColoringDiagram
  \caption{Schematic view of the colored sparse-recovery logic used by the
  local SFD Hessian path: a sparsity pattern, a distance-2 coloring, and the
  corresponding grouped probes used for recovery.}
  \label{fig:globalization-coloring}
\end{figure}

\begin{table}[H]
  \centering
  \begin{tabular}{@{}l c c c >{\RaggedRight\arraybackslash}p{0.38\textwidth}@{}}
\toprule
Family & FEniCS & pure JAX & JAX+PETSc & Higher-order / advanced AD \\
\midrule
$p$-Laplace & yes & yes & yes & element AD and local colored SFD \\
Ginzburg--Landau & yes & no & yes & element AD and local colored SFD \\
Hyperelasticity & yes & yes & yes & trust-region element AD \\
Plasticity2D & no & no & yes & repository scalarized potential and same-mesh PMG \\
Plasticity3D & no & no & yes & constitutive AD, element AD, and same-mesh PMG \\
Topology optimization & no & yes & yes & distributed design updates and PETSc mechanics \\
\bottomrule
\end{tabular}

  \caption{Implementation capability matrix across the maintained benchmark
  families. ``Higher-order / advanced AD'' summarizes the derivative choices
  that are scientifically interesting in each family.}
  \label{tab:capability-matrix}
\end{table}

\Cref{tab:capability-matrix} summarizes the implemented surface. The most
important structural point is that the maintained production story is always
the \jaxpetsc{} column, even when a family also has \fenics{} or \purejax{}
references.
